bundles over $TM\setminus\{0\}$, we will omit mentioning it
explicitly. For example, we will denote by $T^*_v$ the vector
bundle of semi-basic $1$-forms $T^*_v(TM\setminus\{0\})$, which is
a subbundle of $T^*(TM\setminus\{0\})$. We will denote by
$\Lambda^kT^*_v$ the vector bundle of semi-basic $k$-forms on
$TM\setminus\{0\}$, and by
$\Lambda^k_v=Sec\left(\Lambda^kT^*_v\right)$ the
$C^{\infty}(TM\setminus\{0\})$-module of sections $\Lambda^k_v(TM\setminus\{0\})$. By $S^kT^*$ we denote the vector
bundle of symmetric tensors of $(0,k)$-type on $TM\setminus\{0\}$.

The partial dif\/ferential operator $P_1$ induces a morphism of vector
bundles
\[
 p^0(P_1): \ J^1T^*_v \to F_1:=T^*_v \oplus \Lambda^2 T^*_v \oplus
\Lambda^2 T^*_v.
\] Together with this morphism we will consider the $l$th order jet prolongations $p^l(P_1):
J^{l+1}T^*_v \to J^{l}F_1$, for $l\geq 1$.

Locally, for a semi-basic $1$-form $\theta=\theta_idx^i\in
\Lambda^1_v$, we have
\begin{gather*}
\mathcal{L}_{\mathbb C}\theta=\frac{\partial \theta_i}{\partial
  y^j}y^jdx^i, \qquad d_J\theta = \frac{1}{2}\left(\frac{\partial
\theta_i}{\partial y^j}-\frac{\partial \theta_j}{\partial
y^i}\right)dx^j\wedge dx^i, \\ d_h\theta = \frac{1}{2}\left(\frac{\delta \theta_i}{\delta x^j}-
\frac{\delta \theta_j}{\delta x^i} \right) dx^j\wedge dx^i.
\end{gather*}
Therefore, the vector bundle morphism $p^0(P_1)$ can be expressed as
follows
\begin{gather*}
p^0(P_1)(j^1\theta)= \left(\frac{\partial
\theta_i}{\partial y^j}y^j dx^i, \frac{1}{2}\left(\frac{\partial
\theta_i}{\partial y^j}-\frac{\partial \theta_j}{\partial
y^i}\right) dx^j\wedge dx^i, \frac{1}{2}\left(\frac{\delta \theta_i}{\delta x^j}-
\frac{\delta \theta_j}{\delta x^i}\right) dx^j\wedge dx^i \right).
\end{gather*} The symbol
of $P_1$ is the vector bundle morphism $\sigma^1(P_1): T^*\otimes
T^*_v \to F_1$, def\/ined by the highest order terms of $p^0(P_1)$.
Since all terms that def\/ine $p^0(P_1)$ are f\/irst-order terms,
it follows that
\begin{gather} \sigma^1(P_1)A=\left(
\sigma^1\left(\mathcal{L}_{\mathbb C}\right)A=i_{\mathbb{C}}A,
\sigma^1\left(d_J\right)A=\tau_JA,
\sigma^1\left(d_h\right)A=\tau_hA\right). \label{sigma1p1}
\end{gather}
In view of formula \eqref{taull}, the three components of the vector
bundle morphism $\sigma^1(P_1)$ are given~by:
\begin{gather*}
\left( \sigma^1\left(\mathcal{L}_{\mathbb C}\right)A\right) (X)
 = \left( i_{\mathbb{C}}A\right)(X)=A(\mathbb{C}, X); \\
\left(\sigma^1\left(d_J\right)A\right)(X,Y) = \left(\tau_JA\right)(X,Y)
= A(JX,Y)-A(JY, X); \\
\left(\sigma^1\left(d_h\right)A\right)(X,Y) = \left(\tau_hA\right)(X,Y)
= A(hX,Y)-A(hY, X),
\end{gather*} for $X$, $Y$ vector f\/ields on $TM\setminus\{0\}$.
Note that for $A\in T^*\otimes T^*_v$, $i_{\mathbb{C}}A$, $\tau_JA$,
$\tau_hA$ are semi-basic forms and hence the symbol $\sigma^1(P_1)$ is well
def\/ined.

The f\/irst-order prolongation of the symbol of $P_1$ is the vector
bundle morphism $\sigma^2(P_1): S^2T^*\otimes T^*_v \to T^*\otimes
F_1$ that satisf\/ies
$i_X\left(\sigma^2(P_1)B\right)=\sigma^1(P_1)(i_XB)$ for all $B\in
S^2T^*\otimes T^*_v$ and all $X\in {\mathfrak
X}(TM\setminus\{0\})$. Therefore, we obtain
\begin{gather*}
\sigma^2(P_1)B=\left( \sigma^2\left(\mathcal{L}_{\mathbb
C}\right)B, \sigma^2\left(d_J\right)B,
\sigma^2\left(d_h\right)B\right),
\end{gather*} where  for $X, Y, Z$ vector f\/ields on
$TM\setminus\{0\}$ we have:
\begin{gather} \nonumber \left(\sigma^2\left(\mathcal{L}_{\mathbb
C}\right)B\right)(X,Y)   =   B(X, \mathbb{C}, Y), \\
\label{sigma2p1} \left(\sigma^2\left(d_J\right)B\right)(X,Y,Z)   =
B(X, JY,Z)-B(X,JZ,Y), \\
\nonumber \left(\sigma^2\left(d_h\right)B\right)(X,Y,Z)   =
B(X,hY,Z)-B(X,hZ, Y).
\end{gather}
